\documentclass{article}
\usepackage[utf8]{inputenc}
\usepackage[T1]{fontenc}
\usepackage[polish]{babel}
\usepackage{pgfplots}

\pgfplotsset{compat=1.18} 

\begin{document}

% --- POCZĄTEK ROZDZIAŁU DO SKOPIOWANIA ---

\section{Analiza potrzeb użytkowników i identyfikacja problemów}

\subsection{Metodologia badania pilotażowego}
Aby projektowana aplikacja odpowiadała na rzeczywiste problemy studentów Wydziału MiNI, zrezygnowano z opierania logiki nawigacyjnej wyłącznie na statycznych planach ewakuacyjnych budynku. Zamiast tego przeprowadzono badanie pilotażowe na grupie docelowej ($N=100$ studentów pierwszego roku). Głównym celem zwiadu badawczego była identyfikacja tzw. krytycznych węzłów nawigacyjnych, czyli miejsc, w których architektura gmachu najczęściej wprowadza użytkowników w błąd. 

\subsection{Wyniki: Krytyczne węzły nawigacyjne (Pain points)}

Zgromadzone dane ilościowe zostały poddane agregacji i przedstawione na poniższym wykresie.

\begin{figure}[htpb]
    \centering
    \begin{tikzpicture}
        \begin{axis}[
            xbar,                   
            y dir=reverse,          
            xlabel={Odsetek studentów wskazujących salę (\%)},
            ylabel={Numer sali dydaktycznej},
            symbolic y coords={313, 107, 303, 216, 215}, 
            ytick=data,             
            nodes near coords,       
            nodes near coords align={horizontal},
            bar width=18pt,          
            width=12cm,             
            height=7cm,              
            xmin=0, xmax=45,        
            title={Wyniki ankiety: Gdzie studenci gubią się najczęściej?},
            grid=major,             
            major grid style={dotted, gray}
        ]
        
        \addplot[fill=red!70!black, draw=black] coordinates {
            (35,313)    
            (22,107)    
            (18,303)    
            (10,216)    
            (8,215)     
        };
        \end{axis}
    \end{tikzpicture}
    \caption{Rozkład odpowiedzi na pytanie o najtrudniejszą do zlokalizowania salę ($N=100$).}
    \label{fig:trudne_sale}
\end{figure}

Wykres z ankiety (Rysunek \ref{fig:trudne_sale}) jasno pokazuje jedną rzecz: problemem na MiNI nie jest wcale wielkość gmachu, tylko jego nieintuicyjne zakamarki. Skoro na pierwszym miejscu ląduje sala 313, a zaraz za nią 107 z zupełnie innego piętra, to widać, że studenci nie gubią się ogólnie, tylko w bardzo konkretnych, dziwnych miejscach. Numery pokoi po prostu nie zawsze układają się tak, jak podpowiadałaby logika korytarzy.

\subsection{Wnioski dla logiki i architektury systemu}
Wniosek dla naszego projektu jest prosty -- zwykła, płaska mapa z zaznaczoną trasą to za mało. Skoro sam układ korytarzy potrafi zmylić, nasza aplikacja musi być o krok przed użytkownikiem. Jeśli student wyznacza trasę do tej nieszczęsnej sali 313, algorytm nie powinien tylko rysować suchej wektorowej kreski na ekranie. 

System musi natychmiast wyłapać, że węzeł docelowy posiada wysoką ,,wagę trudności'' i dorzucić dodatkową podpowiedź. Może to być zdjęcie kłopotliwego rozwidlenia wyświetlane w interfejsie wariantu trasy, albo szybki komunikat tekstowy w stylu: ,,za automatami z kawą skręć w lewo''. Krótko mówiąc: nasza nawigacja ma celowo łatać błędy w UX (User Experience) samego budynku i prowadzić za rękę dokładnie tam, gdzie statystyki pokazują, że najłatwiej się zgubić.

% --- KONIEC ROZDZIAŁU ---

\end{document}